\cleardoublepage
\thispagestyle{empty}
\begin{center}


{\Large\textbf{Elektroniklabor– Theorie und Anwendungen}}\\[1.2em]
{\large Dipl.-Ing.\,(FH)\,Christian Weilharter}\\[1.2em]
\textcopyright~2025, Christian Weilharter, Traunstein\\[2em]
\end{center}
\begin{flushleft}
	\begin{tabular}{@{}l l}
	
	
		\textbf{ISBN: (Print)} & 978-3-912302-+++\\[0.5em]
		%	\textbf{ISBN: (E-Book)} & 978-3-912302-01-1 \\[0.5em]
		\textbf{Auflage:} & 1.~Auflage 2026 \\[0.5em]
		\textbf{Satz:} & \LaTeX \\[0.5em]
		\textbf{Verlag:} & Christian Weilharter \\[0.5em]
		\textbf{Druck:} & Amazon KDP (Print-on-Demand) \\[0.5em]

		
		\textbf{Kontakt:} & \href{mailto:info@mathandphysics.de}{info@mathandphysics.de}\\[0.5em]
		\textbf{Web:} & \href{https://www.mathandphysics.de}{www.mathandphysics.de}\\
		

	\end{tabular}
\end{flushleft}

\vspace{2em}
\noindent
Alle Rechte vorbehalten. Kein Teil dieses Buches darf ohne schriftliche Genehmigung des Autors 
in irgendeiner Form reproduziert, gespeichert oder übertragen werden, 
weder elektronisch, mechanisch, durch Fotokopien, Aufnahmen noch auf andere Weise.
\begin{center}\small Printed in Germany\end{center}

\cleardoublepage
\chapter*{Vorwort}
\addcontentsline{toc}{chapter}{Vorwort}

Das Elektron gehört zu den unscheinbarsten und zugleich folgenreichsten Objekten der modernen Physik. 







\begin{flushright}
	\textit{Dipl.-Ing.(FH) Christian Weilharter} \\
	\vspace{0.5em}
	[Traunstein, 2025]
\end{flushright}

\newpage
\noindent
\bigskip
\noindent
\textbf{Lesehinweis}

Dieses Buch folgt einem durchgehenden roten Faden.
Begriffe werden eingeführt,
an Beispielen entfaltet
und am Ende der Kapitel in ihrer Bedeutung verdichtet.
Die Boxen im Text markieren dabei zentrale Punkte
wie Definitionen, historische Einordnungen, Denkfehler oder Kernaussagen.
Sie dienen der Orientierung,
ersetzen aber nicht den Fließtext,
in dem die eigentliche Argumentation entwickelt wird.

\medskip
\noindent
Zur besseren Übersicht sind die Boxentypen farblich unterschieden:
Hinweisboxen geben Orientierung,
Historienboxen ordnen Begriffe und Entwicklungen geschichtlich ein,
Matheboxen bündeln formale Inhalte,
und Didaktikboxen heben gedankliche Schlüsselpunkte hervor.
